\section{\hesp{} Design}
\label{sec:approach}

\hesp{} sits between a local LLM and the probe catalogue (\cref{fig:pipeline}). The model proposes: at each step it may suggest a probe, propose a verdict, or abstain, and it parses raw logs when a probe returns text. The controller decides: it chooses which probe runs, which verdict is accepted, and when the investigation ends. Every step is written to a journal before its observation arrives and is audited afterwards.

We walk through the running example. The alert leaves nine candidate causes at probability 0.111 each. The controller ranks the probes by expected information gain per cost and runs \texttt{source\_ips}, although the model proposed \texttt{auth\_log} (1.74 against 0.51 bits per cost unit); the result, \textit{normal}, leaves four causes at 0.238. \texttt{access\_pattern} leaves three at 0.327: the exposed admin panel, a vulnerable component, and DNS command-and-control. \texttt{admin\_exposure} returns \textit{exposed, no auth}, which moves the admin panel to 0.966. Before the next model call, the controller sees that the leader exceeds the acceptance threshold and that the last observation singles it out, and it ends the episode with that verdict after three probes and 5.1\,s.

\begin{figure*}[t]
  \centering
  \includegraphics[width=\textwidth]{fig_pipeline.pdf}
  \caption{\hesp{} on the running example (Llama-3.1-8B; true cause: an exposed admin panel). \protect\stepnum{1} The ledger keeps a posterior over the candidate causes, updated by Bayes' rule from frozen probe-outcome tables. \protect\stepnum{2} The controller runs the legal probe with the highest expected information gain per cost; the model's proposal (here \texttt{auth\_log}, three times) is advisory. \protect\stepnum{3} Once the leader's posterior reaches 0.8 and a current observation singles it out, the controller concludes and cites that evidence. The verifier is used for evaluation only. All values are taken from the archived journal.}
  \label{fig:pipeline}
\end{figure*}

\subsection{Evidence Ledger}
\label{sec:ledger}
\label{sec:tables}
\emph{Problem.} The controller needs a model of what each probe would return under each cause; the LLM cannot provide one (tables elicited from the planner leave resolution at 0.069 to 0.375, \S\ref{sec:policy}).

\emph{Mechanism.} Before deployment, a script runs every probe once in each of $k$ development episodes of every cause and counts the outcomes into a table
\[\widehat P(o \mid h, a) = \frac{q_{h,a,o}}{\sum_{o'} q_{h,a,o'}},\quad q_{h,a,o} = (1-\varepsilon)\,\frac{n_{h,a,o}}{n_{h,a}} + \frac{\varepsilon}{|O_a|},\]
where $n_{h,a,o}$ counts valid observations of outcome $o$ for probe $a$ under cause $h$, $O_a$ is the probe's outcome vocabulary, and $\varepsilon = 0.01$. A cause without development cases, such as \textit{other}, gets a uniform row. The tables are hashed and frozen before use. For each alert, the ledger starts from a uniform prior over the candidate causes and applies Bayes' rule, $p_t(h) \propto p_{t-1}(h)\,\widehat P(o_t \mid h, a_t)$, after every valid observation, recording whether the observation raised or lowered each cause's score. Transient failures (e.g., HTTP 503) are journaled but are not evidence. When the environment reports that the incident changed, the scores reset to the prior and earlier observations stop counting as current.

\emph{Alternatives.} Asking the model for likelihoods removes the need for development data but fails, as noted above. In our settings, five development episodes per cause already give tables that resolve every case (\S\ref{sec:policy}); a deployment would count resolved tickets or replayed incidents (\S\ref{sec:discussion}).

\subsection{Probe Selection}
\label{sec:selection}
For every legal probe, i.e., one that has not run in the current state, is affordable, and has its prerequisites met, the controller computes the expected information gain
\begin{align*}
\mathrm{EIG}(a) &= H(p) - \textstyle\sum_{o} P(o \mid a)\, H\big(p(\cdot \mid o, a)\big),\\
P(o \mid a) &= \textstyle\sum_h p(h)\, \widehat P(o \mid h, a),
\end{align*}
and runs the probe with the highest $\mathrm{EIG}(a)/c(a)$, where $c(a)$ is its cost. The criterion is the classical one~\citep{dekleer1987,golovin2010}; what matters here is that the model's proposal no longer decides which probe runs. The model receives the task, the catalogue, the observations, the ledger, and feedback on rejected proposals (Appendix~\ref{app:prompt}), and returns one JSON decision.

\subsection{Evidence Rule and Stopping}
\label{sec:finish}
\emph{Finish guard.} A verdict proposed by the model is accepted only if the claimed cause has a current score of at least $\tau = 0.8$ and at least one cited observation is current and supports it in the ledger. A rejected verdict is returned to the model as feedback. Text in a log never enters the ledger, so text alone cannot satisfy the guard.

\emph{Controller stop.} Before each model call, and when a budget runs out, the controller applies the same rule to the current leader and, if it holds, concludes on its own, citing up to three supporting observations. The \emph{confirmation} variant additionally requires one observation that \emph{specifically} supports the leader, i.e., whose outcome is at least ten times more likely under the leader than under any other named cause. It uses the tables only and forbids conclusions reached purely by elimination.

\emph{Recovery.} A model can stall a guarded controller by proposing the same verdict just below the threshold. After two consecutive rejected verdicts, the controller runs its own top-ranked probe instead of asking the model again.

\subsection{Source-Diverse Corroboration}
\label{sec:corroboration}
Once the controller decides, an attacker forges evidence instead of text, and a single forged observation that singles out a benign cause suffices, because the controller acts on its most informative observation first. \hesp{} therefore requires a benign verdict, whether proposed by the model or reached by the controller stop, to rest on current observations from at least two different \emph{source groups} that each specifically support it. A source group is the upstream data source behind a probe, declared per deployment; two probes that read the same reputation feed form one group. If the rule is not met, the controller keeps probing, and if the evidence stays inconsistent, the case is escalated.

\subsection{Reader Trust}
\label{sec:readertrust}
When a probe returns log text, a parser maps it to the probe's outcome vocabulary. A rule parser written for the documented format is safe but reads nothing after a format change; the model reads any format but can be persuaded. \hesp{} tags each observation with the parser that produced it. Under the \emph{reader-trust} rule, an observation parsed by the model updates the posterior and may confirm an actionable cause, but it never counts as specific support for a benign verdict, in the controller stop, in corroboration, and in the finish guard alike. The asymmetry follows the attacker's goal: a forged reading that confirms an attack costs an unnecessary escalation, whereas one that confirms a benign cause closes the case. The rule is a provenance label in the spirit of capabilities and information-flow labels~\citep{camel,fides}, specialized to one decision.

\subsection{Security Analysis}
\label{sec:analysis}
Whatever the model outputs, three properties hold: every executed probe is in the read-only catalogue and affordable; every accepted verdict cites a current observation that raised the claimed cause's score (G3); and every prediction is journaled before the observation it predicts, which the audit checks. Against A1, the finish guard and the controller stop suffice, because injected text never becomes ledger evidence. Against A2, neither helps: a forged source produces genuine-looking evidence. Source-diverse corroboration restores G2 as long as the attacker controls at most one source group. Against A3, corroboration no longer helps, because the attacker's text arrives with every probe's response and can forge the reading of every source group. Reader trust restores G1 as long as a trusted parser exists, at the cost of escalating honest benign cases whenever only the model can read the logs. None of these rules makes a verdict correct; correctness depends on the tables and on the true cause being among the candidates, which \S\ref{sec:policy} evaluates.

\subsection{Implementation}
\label{sec:implementation}
\hesp{} is about 2{,}900 lines of Python that use only the standard library. One JSON decision protocol serves vLLM and Ollama endpoints and the scripted planners used as LLM-free references. Each episode runs against its own loopback HTTP sandbox, so probes are real HTTP requests with transient failures and state changes. A write-ahead journal records every request, decision, prediction, observation, and ledger update, and an audit replays each journal and checks event order, provenance, budgets, citations, and every controller verdict. 201 unit tests cover the ledger, selectors, guard, stop, corroboration, reader trust, and the isolation of the table estimator from the environment's generator.
